\section{Design}
\label{sec:design}

This section is the core of the report. It first explains why the architecture proposed
in the first progress journal was replaced (\S\ref{sec:design-pivot}), then states the
design principles that the replacement follows (\S\ref{sec:design-principles}), and
finally presents each mechanism with its rationale
(\S\ref{sec:design-coalescing}--\S\ref{sec:design-priority}).

\subsection{Where should the meter live?}
\label{sec:design-pivot}

\paragraph{The Journal-1 proposal.}
The first progress journal framed the project as ``Postage: paid attention'' and proposed
a standalone TypeScript service in the mould of Tack. A sender wishing to reach a receiver
would call the service, receive an x402 \code{402 Payment Required} with a price, pay in
USDC on Taiko Alethia, and retry; the service would then relay or release the message.
The proposal had real merits: it reused a payment flow the partner already operated, it
put Taiko in the per-message path, and it could be built without modifying TAP.

\paragraph{Why it was set aside.}
Working through the design against the codebase exposed four objections, any one of
which would have been sufficient.

\begin{enumerate}
  \item \textbf{It meters the wrong resource in the wrong place.} Tack prices bytes stored
        at Tack. The resource this project is about is the \emph{receiver's} LLM attention,
        and only the receiver's daemon knows what a message actually cost: whether it was
        rendered at all, into how many lines, whether it coalesced with earlier lines, whether
        it triggered a wake-up, whether it was suppressed past the cap. A service in front
        of the receiver can count messages or bytes; it cannot count attention. Pricing
        attention from outside the receiver would be pricing a proxy.
  \item \textbf{It duplicates a trust layer TAP already has.} TAP admits a sender only
        after an ERC-8004 identity has been resolved and a signed invite exchanged, and it
        already records per-peer grants. A standalone service would need its own account
        model, its own authentication of both parties and its own mapping from payer to
        recipient, all of which would have to be kept consistent with the trust store inside
        each daemon. It would also be a new party that both peers must trust with
        message-flow metadata, contradicting TAP's local-first design and NFR4.
  \item \textbf{It puts chain settlement in the hot path.} A per-message x402 payment means a
        transaction, or at least a facilitator round-trip, per message. The upstream TAP
        maintainers' research note for the messaging layer states the opposite requirement: ``Relays only
        need to verify the proof and decrement balance. No chain RPC is needed in the hot
        path''~\cite{tap-research-note}. The economically sound shape is one payment
        converted into many cheap spends.
  \item \textbf{It cannot compose with grants.} A friend who has been granted
        \code{message/send} should never pay; a stranger should. Only the receiver knows who
        holds a grant. Enforcing payment outside the receiver would either charge friends or
        require the receiver to publish its grant table to the service.
\end{enumerate}

\paragraph{The decision.}
Postage was therefore built \emph{into} TAP, at the receiver, as an extension of the
daemon that already renders notifications. Figure~\ref{fig:pivot} contrasts the two
shapes. What was kept from the x402 model is its most useful idea: an unpaid request is
answered not with silence but with a machine-readable quote. In the implemented system
that quote travels as JSON-RPC error \code{-32050} (\code{ATTENTION\_PAYMENT\_REQUIRED})
over XMTP, and the CLI and daemon surface it to local callers as HTTP \code{402}. What was
kept from the partner relationship is the settlement leg: a postage top-up is an on-chain
USDC payment to the receiver's agent address on whichever CAIP-2 chain the receiver
registered on, so an agent registered on Taiko is paid on Taiko. This was a decision taken
with the project's full time budget in hand; it made the resulting system smaller, not
larger, but the reason for it is correctness, not scope.

\begin{figure}[htbp]
\centering
\begin{tikzpicture}[
  font=\small,
  box/.style={draw, rounded corners=2pt, minimum height=9mm, minimum width=26mm, align=center, inner sep=3pt},
  ext/.style={box, fill=red!8, draw=red!60!black},
  tap/.style={box, fill=blue!6, draw=blue!50!black},
  chain/.style={box, fill=green!8, draw=green!50!black},
  arr/.style={-{Stealth[length=2mm]}, thick},
  lbl/.style={font=\scriptsize, fill=white, inner sep=1.5pt},
]
% ---------------- (a) Journal-1 plan ----------------
\node[font=\bfseries\small, anchor=west] at (-6.4,1.9) {(a) Journal-1 plan: standalone x402 service between the peers};
\node[tap] (s1) at (-4.6,0.6) {Sender\\daemon};
\node[ext] (svc) at (0,0.6) {Tack-like\\x402 service};
\node[tap] (r1) at (4.6,0.6) {Receiver\\daemon};
\node[chain] (c1) at (0,-1.6) {Taiko Alethia\\(USDC)};
\draw[arr] (s1) -- node[lbl,above]{send} (svc);
\draw[arr] (svc) -- node[lbl,above]{release} (r1);
\draw[arr,dashed] (svc) -- node[lbl,right]{402 quote / verify, per message} (c1);
\draw[arr,dashed] (s1.south) |- node[lbl,below,pos=0.75]{pay per message} (c1.west);
\node[font=\scriptsize, align=left, anchor=west, text width=52mm] at (1.6,-1.6)
  {Meter counts messages or bytes at a third party; grants are invisible to it; one settlement per message.};
% ---------------- (b) Implemented ----------------
\begin{scope}[yshift=-5.6cm]
\node[font=\bfseries\small, anchor=west] at (-6.4,1.9) {(b) Implemented: TAP-native postage at the receiver};
\node[tap] (s2) at (-4.6,0.6) {Sender\\daemon};
\node[tap, minimum width=46mm] (r2) at (2.0,0.6) {Receiver daemon\\\scriptsize gate $\to$ ledger $\to$ render};
\node[chain] (c2) at (-1.2,-1.6) {Receiver's chain\\(Base or Taiko)};
\draw[arr] ([yshift=2mm]s2.east) -- node[lbl,above]{stamped \code{message/send}} ([yshift=2mm]r2.west);
\draw[arr,dashed] ([yshift=-2mm]r2.west) -- node[lbl,below]{\code{-32050} quote} ([yshift=-2mm]s2.east);
\draw[arr,dashed] (s2.south) |- node[lbl,below,pos=0.75]{top-up once} (c2.west);
\node[font=\scriptsize, align=left, anchor=west, text width=52mm] at (1.6,-1.6)
  {Meter counts rendered lines and tokens at the receiver; grants are free stamps; one settlement per credit.};
\end{scope}
\end{tikzpicture}
\caption{Placement of the attention meter. (a)~The first progress journal proposed a
standalone x402 service between the peers. (b)~The implemented design places metering,
pricing and postage inside the receiver's daemon, where attention is actually spent and
where identity and grants already live; the chain is touched once per credit, not once per
message.}
\label{fig:pivot}
\end{figure}

\subsection{Design principles}
\label{sec:design-principles}

The principles below were written down before the implementation began and were used to
settle the many smaller decisions that follow. Each is stated with the requirement it
serves.

\begin{enumerate}[label=P\arabic*.]
  \item \textbf{Meter where the cost lands.} Attention is accounted at the moment a line is
        rendered into a host context, from the same render plan the host uses, so the ledger
        can never disagree with what the agent saw (FR3, FR11).
  \item \textbf{Compress before you charge.} Coalescing and escalation-first rendering come
        first, because they eliminate most of the waste at no cost to anyone and because
        pricing without a unit of measurement is guesswork. Only after the ledger exists is
        it meaningful to publish a price (FR1--FR4).
  \item \textbf{Escalations outrank everything.} No mechanism---cap, coalescing, quota---may
        hide an escalation. Paid priority stamps and pending approvals are always rendered
        and never folded (FR2, FR9, FR10).
  \item \textbf{Never lose a message; fail open.} Metering may fold notifications into
        summaries; it may not drop messages. If a quota or ledger lookup fails, the batch
        passes through unmetered (NFR3).
  \item \textbf{Social trust precedes economic trust.} A grant is a free stamp. Postage
        exists for un-granted peers, and for granted peers who want more than the grant
        gives (a wake-up, or lines beyond the quota) (FR6, FR9, FR10).
  \item \textbf{Pay once, spend many.} One on-chain payment opens a credit; stamps are
        debited off-chain under a local mutex. Nothing in the per-message path touches the
        chain (FR7).
  \item \textbf{Opt in, and be machine-readable when you say no.} Enforcement is off by
        default; when it rejects, the rejection carries everything the sender needs to fix
        the situation without a human reading logs (FR5, NFR1).
  \item \textbf{Idempotent under crash and redelivery.} Every step that moves value or
        state is keyed so that a retry after a crash converges rather than double-charges
        (NFR2).
\end{enumerate}

\subsection{Pipeline overview}

Figure~\ref{fig:pipeline} shows where each mechanism sits in the inbound path. The gate
(\S\ref{sec:design-pricing}, \S\ref{sec:design-postage}) runs before anything is logged;
coalescing (\S\ref{sec:design-coalescing}) runs at enqueue; quota folding
(\S\ref{sec:design-priority}), accounting (\S\ref{sec:design-ledger}) and the
escalation-first render plan (\S\ref{sec:design-coalescing}) run at drain, inside the
daemon, so that every host receives the same batch.

\begin{figure}[htbp]
\centering
\resizebox{\textwidth}{!}{%
\begin{tikzpicture}[
  font=\small,
  node distance=7mm and 16mm,
  stage/.style={draw, rounded corners=2pt, minimum height=10mm, minimum width=27mm, align=center, inner sep=3pt, fill=blue!5, draw=blue!50!black},
  new/.style={stage, fill=orange!12, draw=orange!70!black},
  arr/.style={-{Stealth[length=2mm]}, thick},
  lbl/.style={font=\scriptsize, fill=white, inner sep=1.5pt},
]
\node[stage] (xmtp) {XMTP\\transport};
\node[new, right=of xmtp, minimum width=34mm] (gate) {Attention gate\\\scriptsize grant? stamp? else \code{-32050}};
\node[stage, right=of gate] (log) {Conversation\\log + journal};
\node[stage, right=of log] (evt) {Typed event\\\scriptsize \code{message.received}};
\node[new, below=16mm of evt] (cls) {Classifier\\\scriptsize priority $\Rightarrow$ escalation};
\node[new, left=of cls] (q) {Coalescing\\queue\\\scriptsize merge by key};
\node[new, left=of q, minimum width=34mm] (drain) {Drain\\\scriptsize quota fold $\to$ account};
\node[new, left=of drain] (plan) {Render plan\\\scriptsize escalations first, cap 20};
\node[stage, below=14mm of plan, xshift=-8mm, minimum width=48mm] (host) {Host prompt\\\scriptsize OpenClaw / Hermes: \code{[TAP Notifications]}};
\node[new, below=14mm of drain, xshift=10mm, minimum width=34mm] (ledger) {Attention ledger\\\scriptsize per peer, per UTC day};
\draw[arr] (xmtp) -- (gate);
\draw[arr] (gate) -- node[lbl,above]{free or paid} (log);
\draw[arr] (log) -- (evt);
\draw[arr] (evt) -- (cls);
\draw[arr] (cls) -- (q);
\draw[arr] (q) -- (drain);
\draw[arr] (drain) -- (plan);
\draw[arr] (plan.south) -- (host.north -| plan.south);
\draw[arr] (drain.south) -- (ledger.north -| drain.south);
\draw[arr, dashed, red!70!black] (gate.south) -- ++(0,-0.75) node[lbl, below, text=red!70!black]{rejected: not logged, not rendered, not billed};
\end{tikzpicture}%
}
\caption{Inbound path in the receiving daemon. Orange stages were added by this project.
A rejected message never reaches the log, the queue or the ledger. Everything from the
queue to the render plan runs inside \code{tapd}, so the CLI, OpenClaw and Hermes render
one canonical batch.}
\label{fig:pipeline}
\end{figure}

\subsection{Notification coalescing and escalation-first rendering}
\label{sec:design-coalescing}

\paragraph{Problem restated.}
The legacy queue appended one notification per event and the renderer emitted the first
twenty in arrival order. Fifty chat messages from three peers followed by one transfer
request produced twenty interleaved chat lines and a footer ``31 more TAP notifications
omitted''---with the approval request inside the 31.

\paragraph{Coalescing at enqueue.}
Each notification may carry a \code{coalesceKey} and a \code{coalesceStrategy}. When a
new notification arrives whose key is already buffered, it is merged \emph{in place}: the
entry keeps its buffer slot and identifier, its content becomes the newest event's
content, and, under the default \code{count} strategy, a \code{count} accumulates. The
alternative \code{replace} strategy swaps content without counting; it is used for state
transitions such as an action going from pending to completed, which is one thing changing
rather than many things repeating. Chat from one connection is keyed
\code{msg:<connectionId>}, so a flood from Alice becomes a single line
``\code{New message from Alice: \ldots (x334)}''. The design choice to keep the
\emph{first-arrival} slot but the \emph{newest} content is deliberate: the slot preserves
fairness of ordering across peers, while the content shows the operator the most recent
state of the conversation.

Coalescing happens at enqueue rather than at drain for two reasons. First, the queue is
bounded (1{,}000 entries, oldest dropped) as a safety measure against a host that never
drains; merging at enqueue means a flood from one peer occupies one slot and cannot evict
other peers' notifications. Second, a per-key index kept in step with the buffer makes
the merge $O(1)$; the index is cleared on drain and pruned on eviction so that a stale
mapping can never merge new events into an entry that is no longer buffered.

\paragraph{Escalation-first render plan.}
Rendering was refactored into a pure function that returns a \emph{plan}
(\code{planNotificationRender}): the lines that will be rendered, the entries omitted
past the cap, and a footer. The plan drops empty one-liners, then stable-partitions escalations ahead of all
other types (preserving arrival order within each class), caps at twenty, and summarises
the omitted tail \emph{grouped by peer and type with event counts}
(``\code{omitted: 334 messages from Alice, 2 auto-replies}'') rather than as an opaque
number. Because the plan is a value rather than a side effect, the OpenClaw host renders
exactly it, the Hermes Python plugin mirrors it, and the accounting step
(\S\ref{sec:design-ledger}) bills from it. This single-source-of-truth property is what
satisfies FR11.

\paragraph{Priority messages are never coalesced.}
A message whose stamp paid the priority tier (\S\ref{sec:design-priority}) is classified
as an escalation with a 400-character excerpt instead of 80 and carries no coalesce key.
A paid wake-up must remain individually visible; collapsing two of them into one line
would silently refund nothing and hide one.

\subsection{The attention ledger}
\label{sec:design-ledger}

The ledger answers one question: \emph{what has each peer cost me?} It is a file
(\path{<dataDir>/attention-ledger.json}) of UTC-day buckets, each holding per-peer
statistics and a total. The unit of account is defined by the render plan: for every
rendered line the ledger adds one to \code{notificationsRendered} and
$\lceil \text{chars}/4 \rceil$ to \code{tokensInjected}; escalations increment
\code{escalations}; omitted entries add their event counts to
\code{overflowSuppressed}. Peers are keyed \code{<chain>\#<agentId>}, the same key the
postage ledger uses so the two can be joined; block overhead and events with no peer (a
pending-action escalation, for instance) are booked to an \code{unattributed} row.

Three properties are intentional. The ledger records \emph{rendered} lines, not received
messages, so a peer whose flood coalesced into one line is billed one line---the ledger
measures attention actually spent, which is what P1 demands and what makes quota
enforcement (\S\ref{sec:design-priority}) fair. Retention is thirty days with pruning on
write, so the file is bounded. And the ledger is written only by the transport-owning
daemon, under a process-local mutex and the repository's atomic \code{tmp + rename}
pattern, so it inherits the single-writer discipline TAP already imposes on the trust
store. \code{tap attention show} is a pure local read and starts no transport, in keeping
with the ``fat CLI'' rule.

\subsection{Attention pricing and the \code{-32050} quote}
\label{sec:design-pricing}

\paragraph{Publishing prices.}
An agent's price list lives in \code{config.yaml} under \code{attention.pricing} and is
copied into the registration file as \code{trustedAgentProtocol.attention} on the next
\code{tap register update}, with a \code{version}, a \code{currency} (USDC), an optional
CAIP-2 \code{chain}~\cite{caip2} and a map of tiers to decimal amounts. Three tiers are defined:
\code{grantHolder} (documentary---grant holders are free by construction),
\code{standard} (the price a stamp must cover to buy delivery) and \code{priority} (the
price that additionally buys a wake-up). Putting prices in the registration file rather
than inventing a new discovery message reuses the existing, cached
(\raisebox{0.3ex}{\textasciitilde}24\,h) resolution path: a sender learns a peer's prices
with \code{tap identity resolve}, and \code{tap message send --dry-run} previews the cost
and whether a stamp would be attached.

\paragraph{The gate.}
With \code{attention.enforce: true}, an inbound \code{message/send} passes through a gate
before the conversation log. If the sender holds an active \code{message/send} grant,
the message is free (FR6). Otherwise the gate looks for a stamp
(\S\ref{sec:design-postage}); if none is present, or the stamp does not cover the
\code{standard} price, the message is rejected with
\code{AttentionPaymentRequiredError}, which the transport maps to JSON-RPC
\code{-32050} with \code{error.data.attention} holding the quote and
\code{error.data.postage} holding a reason from the closed set
\{\code{missing\_stamp}, \code{unpriced}, \code{unknown\_credit},
\code{seq\_replayed}, \code{below\_price}, \code{insufficient\_credit}\}. The reason set
was chosen so that a sender can decide mechanically whether a top-up would help
(\code{insufficient\_credit}, \code{unknown\_credit}) or whether the stamp itself is wrong.

Two details carry the compatibility requirement. First, if enforcement is on but no
\code{standard} price is published, the gate degrades to ``granted peers only'' rather
than accepting any stamp: an unpriced receiver cannot be bought. Second, the rejection
marks the request journal entry \emph{completed} before throwing, so a later full-history
sync replays the duplicate short-circuit instead of re-running enforcement per replay.

\subsection{Prepaid postage}
\label{sec:design-postage}

Postage is the two-step model from the research note~\cite{tap-research-note}
instantiated between two peers with no facilitator. Figure~\ref{fig:postage} shows the
lifecycle.

\begin{figure}[htbp]
\centering
\begin{tikzpicture}[
  font=\small,
  arr/.style={-{Stealth[length=2mm]}, thick},
  lbl/.style={font=\scriptsize, fill=white, inner sep=1.5pt, align=center},
]
\def\L{0}\def\C{4.6}\def\R{9.2}
\node[font=\bfseries] (sT) at (\L,0) {Sender A};
\node[font=\bfseries] (cT) at (\C,0) {Chain};
\node[font=\bfseries] (rT) at (\R,0) {Receiver B};
\draw[thick, gray] (\L,-0.3) -- (\L,-8.2);
\draw[thick, gray] (\C,-0.3) -- (\C,-8.2);
\draw[thick, gray] (\R,-0.3) -- (\R,-8.2);
% 1 pay
\draw[arr] (\L,-0.8) -- node[lbl,above]{1. USDC transfer to B's agent address (\code{txHash})} (\C,-0.8);
% 2 record held
\node[lbl, anchor=west, draw=gray, rounded corners=1pt] at (\L+0.15,-1.45) {2. record \emph{held} credit \{creditId, txHash\} before any transport};
% 3 topup request
\draw[arr] (\L,-2.2) -- node[lbl,above]{3. \code{action/request postage/topup} \{creditId, amount, chain, txHash\}} (\R,-2.2);
% 4 verify + record issued + sign
\node[lbl, anchor=east, draw=gray, rounded corners=1pt, text width=44mm] at (\R-0.15,-3.15)
  {4. verify payment (host hook); record \emph{issued} credit under mutex: one credit per txHash, $\le$16 live credits per peer; then sign certificate via OWS};
% 5 result
\draw[arr] (\R,-4.1) -- node[lbl,above]{5. \code{action/result} \{status: accepted, certificate\}} (\L,-4.1);
\node[lbl, anchor=west, draw=gray, rounded corners=1pt] at (\L+0.15,-4.7) {6. verify certificate against B's known agent address; mark held credit verified};
% 7 stamped sends
\draw[arr] (\L,-5.5) -- node[lbl,above]{7. \code{message/send} + stamp \{creditId, seq=1, cost=0.001\}} (\R,-5.5);
\node[lbl, anchor=east, draw=gray, rounded corners=1pt] at (\R-0.15,-6.05) {debit: seq $>$ lastSeq, cost $\ge$ standard, remaining $\ge$ cost};
\draw[arr] (\L,-6.7) -- node[lbl,above]{8. \code{message/send} + stamp \{seq=2\}} (\R,-6.7);
\draw[arr] (\L,-7.5) -- node[lbl,above]{9. stamp \{seq=3\}} (\R,-7.5);
\draw[arr, red!70!black] (\R,-7.95) -- node[lbl,below]{\code{-32050}, postage.reason = \code{insufficient\_credit}, remaining = 0} (\L,-7.95);
\end{tikzpicture}
\caption{Postage lifecycle between two connected agents. One on-chain payment (1) opens a
credit; all later steps are off-chain. Steps 7--9 correspond to the teammate's Experiment~2
(Section~\ref{sec:eval-postage}).}
\label{fig:postage}
\end{figure}

\paragraph{Top-up as a TAP action.}
A top-up is an ordinary \code{action/request} of type \code{postage/topup}, handled by a
TAP app. This was the natural fit: the app registry, the request journal, the
result-waiter machinery and the conversation log all apply without special cases, and the
same handlers ship both as a built-in of \code{core} and as the standalone package
\code{@trustedagents/app-postage} for hosts that assemble their own app set. The handler
receives its capabilities through a narrow, host-injected extension
(\code{ctx.extensions.postage}: \code{recordIssued}, \code{issuedFor},
\code{signCredit}, \code{verifyTopup}) rather than through the generic app storage,
precisely so that top-ups and stamp debits are serialised by \emph{one} ledger mutex.

\paragraph{Pay first, then ask.}
The sender pays before it sends the top-up request, and records the \emph{held} credit
locally before any transport action. The ordering is chosen for recoverability: the
payment is irreversible, so whatever fails afterwards, the sender's ledger holds the
\code{creditId}/\code{txHash} pair needed to retry, and the receiver's handler is
idempotent per \code{creditId}---a retry after a timeout returns the original
certificate and never opens a second credit. A definitive rejection by the peer removes
the held record so that it cannot feed auto-stamping.

\paragraph{The credit certificate.}
On acceptance the receiver signs, through its OWS signing provider, an EIP-191 personal
signature~\cite{eip191} over the keccak digest of the ABI-encoded facts
\code{(creditId, issuerChain, issuerAgentId, holderAgentId, amountMicros, txHash)}. The
holder verifies it against the issuer's agent address already recorded in its contact
entry---the same scheme and the same trust anchor that invites use. The certificate is the
holder's portable proof that the credit exists and is bound to a specific payment; it is
not needed on the per-message path.

\paragraph{Stamps and the debit rule.}
Once a verified held credit exists, \code{sendMessage} attaches a stamp
\code{\{creditId, seq, cost\}} to the \code{trustedAgent} metadata of every
\code{message/send} to that peer automatically, drawing on the oldest credit whose
remaining balance covers the cost, and persists the incremented \code{seq} and spend
\emph{before} the stamp is returned---so a crash wastes one sequence value rather than
reusing one. On the receiving side a debit succeeds only if the credit exists and belongs
to this peer, \code{cost} is well-formed and at least the \code{standard} price,
\code{seq} is strictly greater than the credit's \code{lastSeq}, and the remaining balance
covers \code{cost}. The whole check-and-spend runs under the ledger mutex. The strictly
increasing sequence is the replay defence: a captured stamp cannot be presented twice,
and there is no nonce set to store or expire. The one exception is deliberate: a stamp
with \code{seq == lastSeq} \emph{and} the same delivery key as the last debit is accepted
without charge, because it is the crash-window redelivery of the message that already
paid.

\paragraph{Why stamps live in \code{message/send} metadata.}
Carrying the stamp inside the existing message rather than as a separate action keeps the
per-message path to one wire message, and lets a pre-project receiver ignore it: metadata
it does not understand is simply not read (NFR1).

\subsection{Priority stamps and notification quotas}
\label{sec:design-priority}

The final mechanism separates the two prices identified in
Section~\ref{sec:background}---context and wake-ups---and gives grant holders a metered
free tier.

\paragraph{Priority buys the wake-up, the grant buys delivery.}
A stamp whose \code{cost} meets the receiver's published \code{priority} price is booked
at tier \code{priority}. The classifier then emits an \emph{escalation} with a
400-character excerpt, the queue never coalesces it, the render plan places it first,
and OpenClaw's escalation watcher calls \code{requestHeartbeatNow()} so the agent wakes
immediately (Hermes shows it on the next turn). Critically, this applies to grant holders
too: the gate exempts a grant holder from paying for delivery, but a priority-value stamp
from a grant holder is still debited and still buys the wake-up
(\code{chargeOptionalPriorityStamp}). The grant encodes ``I will read what you send'';
the priority stamp encodes ``interrupt me now''---these are different permissions and are
priced separately. Any failure in the optional priority path degrades silently to plain
grant delivery: a grant holder is never rejected over a malformed stamp.

\paragraph{The weekly quota.}
A \code{message/send} grant may carry the constraint \code{notificationsPerWeek: N}. At
drain time, before accounting and before any host renders, the daemon groups the batch's
\code{info} notifications by peer, asks the attention ledger how many lines that peer had
rendered in the trailing seven UTC days, and---if the peer is at or over quota---replaces
the peer's info notifications with one \code{summary} line
(``\code{12 notifications from Alice folded --- weekly notification quota reached
(20/week)}'') whose \code{count} preserves the number of underlying events. The messages
are still delivered and still in the conversation log; only their per-line claim on
attention is withdrawn. Escalations are excluded from folding by construction (P3). When
several active grants disagree, the most permissive wins, because all of them were issued
by this same agent. If the trust store or ledger cannot answer, the batch passes through
unfolded (P4).

Quota enforcement is placed in the drain rather than in the renderer for the same reason
accounting is: it must happen once, in the daemon, so that OpenClaw, Hermes and the
ledger all see the same folded batch. It counts \emph{rendered} lines rather than
received messages so that a peer whose chatter already coalesced is not penalised twice.
Together the two mechanisms make the grant's free tier finite and let the peer buy past it
with postage---which is exactly the relationship P5 asks for.
